\section{三个合理要求不能同时满足}
\label{sec:17-Constrained-Guarantees/03-Fairness-Impossibility/02}

评审会上有人提了一个看起来最省事的方案：既然双方各拿一个判据，那就把两个判据都写进验收标准，谁也不用让步。会议纪要里写下三条：分数必须在两个群体内都校准；两个群体的真阳率必须相等；两个群体的假阳率必须相等。

这份纪要在数学上是不可执行的。只要两个群体的基率不同，且分类器会犯错，这三条就没有任何一个模型能同时满足——不是难以满足，是没有。下面把它证完，代数只有四行。

\subsection{一个把基率、命中率与预测值绑在一起的恒等式}

固定群体 \(A=a\)，记
\[
\begin{aligned}
p_a&=P(Y=1\given A=a),\\
\mathrm{TPR}_a&=P(\hat Y=1\given Y=1,\,A=a),\qquad
\mathrm{FPR}_a=P(\hat Y=1\given Y=0,\,A=a),
\end{aligned}
\]
再记 \(v_a=P(Y=1\given \hat Y=1,A=a)\) 为该群体内的阳性预测值——被判为高风险的人里真的再犯的比例。这四个量都是同一张 \(2\times2\) 表的边缘量，彼此不独立。

把表里两个正例格子的概率写出来，再对第二个坐标求和：
\[
\begin{aligned}
P(\hat Y=1,\,Y=1\given A=a)&=p_a\,\mathrm{TPR}_a,\\
P(\hat Y=1,\,Y=0\given A=a)&=(1-p_a)\,\mathrm{FPR}_a,\\
v_a&=\frac{p_a\,\mathrm{TPR}_a}{p_a\,\mathrm{TPR}_a+(1-p_a)\,\mathrm{FPR}_a},\\
\frac{1-v_a}{v_a}&=\frac{(1-p_a)\,\mathrm{FPR}_a}{p_a\,\mathrm{TPR}_a}.
\end{aligned}
\]
整理最后一式，得到
\[
\mathrm{FPR}_a=\frac{p_a}{1-p_a}\cdot\frac{1-v_a}{v_a}\cdot\mathrm{TPR}_a .
\]

这是一个恒等式，对任何分类器、任何群体都成立，里面没有一条假设。它说的是：基率的几率 \(p_a/(1-p_a)\)、预测值的几率、以及两个命中率，四者当中只有三个可以自由指定，第四个被前三个定死。

\begin{proposition}
设两个群体的基率满足 \(0<p_1<1\)、\(0<p_2<1\) 且 \(p_1\neq p_2\)，设分类器在两个群体上都有 \(\mathrm{TPR}_a>0\) 与 \(\mathrm{FPR}_a>0\)。则以下三条不能同时成立：阳性预测值相等（\(v_1=v_2\)）、真阳率相等（\(\mathrm{TPR}_1=\mathrm{TPR}_2\)）、假阳率相等（\(\mathrm{FPR}_1=\mathrm{FPR}_2\)）。
\end{proposition}

证明骨架只有一步。若三条同时成立，把共同值记为 \(v,\mathrm{TPR},\mathrm{FPR}\)，对 \(a=1,2\) 分别代入上面的恒等式，得
\[
\mathrm{FPR}=\frac{p_1}{1-p_1}\cdot\frac{1-v}{v}\cdot\mathrm{TPR}
=\frac{p_2}{1-p_2}\cdot\frac{1-v}{v}\cdot\mathrm{TPR}.
\]
因为 \(\mathrm{FPR}>0\)，右边两个因子 \((1-v)/v\) 与 \(\mathrm{TPR}\) 都不为零，可以约掉，于是
\[
\frac{p_1}{1-p_1}=\frac{p_2}{1-p_2}.
\]
函数 \(p\mapsto p/(1-p)\) 在 \((0,1)\) 上严格递增，故 \(p_1=p_2\)，与假设矛盾。

\subsection{条件对账：每一条挡住什么}

\textbf{基率不同这一条是必需的。}若 \(p_1=p_2\)，恒等式对两个群体给出同一个方程，三条要求之间不再有冲突——最简单的例子是让分类器完全忽略 \(A\)，此时两个群体的整张 \(2\times2\) 表按比例一致，四个判据全部成立。所以这条命题不是在说``两个群体不可能被同等对待''，它是在说：当两个群体的 \(Y\) 分布本身不同时，这种不同必须在某个判据上显现出来，你只能选它显现在哪里。

\textbf{分类器不完美这一条也是必需的。}若 \(\mathrm{FPR}_a=0\) 且 \(\mathrm{TPR}_a=1\)，则 \(v_a=1\)，恒等式两边都是零，无论基率差多大都成立。完美分类器同时满足全部判据。命题里 \(\mathrm{FPR}_a>0\) 这个条件正是``会犯错''的形式化写法，去掉它结论就塌了。

\textbf{还有两条隐含条件不能省。}\(0<p_a<1\) 保证每个群体内两类都有人，否则 \(\mathrm{TPR}\) 或 \(\mathrm{FPR}\) 无定义；\(v_a\) 有定义要求该群体确实有人被判为正例。这些不是技术性的装饰——一个从不判正例的分类器平凡地满足很多判据，把它排除在外是必要的。

把这两条必需条件放在一起看，结论就清楚了：\textbf{冲突不来自模型质量，来自问题结构}。再多的数据、再强的模型族、再仔细的调参，都不能把 \(p_1\neq p_2\) 变成 \(p_1=p_2\)，也不能把一个有噪声的预测问题变成可完美分离的问题。这与\hyperref[ch:088]{``数据压缩与可估计信息''}里那种上限性质不同：那里的界说的是这份数据里的信息不够，多采样可以推进；这里的限制是几个定义之间的代数矛盾，采样量不进入结论。也与\hyperref[ch:118]{``统计学习理论：训练误差对真风险有多乐观''}不同：那里的界随样本量收紧，这里的等式不随任何东西变化。

\begin{remark}
上面用的是阳性预测值相等，即``预测平等''。若把要求换成分数层面的完全校准，同样的结论仍然成立，只是真阳率与假阳率要换成期望形式 \(\E[S\given Y=y,A=a]\)。这个版本的证明思路一致：把 \(\E[S\given A=a]\) 按 \(Y\) 分解一次，校准条件让 \(\E[Y\given S]\) 可以换成 \(S\)，两式相减即得基率必须相等。分数层面的版本更强，因为它不依赖阈值怎么选。
\end{remark}

\begin{remark}
一个自然的退路是把``相等''放松成``近似相等''。这条退路是关着的：已知的定量版本表明，若三个判据都只差一个小量 \(\varepsilon\)，则两个群体的基率之差也必须被 \(\varepsilon\) 的某个函数控制住。基率差是给定的、不可调的，所以 \(\varepsilon\) 有一个由数据决定的下界。近似不能绕过精确的不可能性，只能把它翻译成一个数值。
\end{remark}

\subsection{一条曲线降下去，另一条就升上来}

把上述结论画出来，形状是这样的：沿着任何一条把权重从``保校准''移向``拉平假阳率''的调整路径走，两个偏差量此消彼长，中间不存在同时归零的点。

\begin{center}
\begin{tikzpicture}[x=1cm,y=1cm,font=\small,>=stealth]
  % ---- 左图：基率不同 ----
  \begin{scope}
    \draw[->,black!60] (0,0) -- (4.7,0);
    \draw[->,black!60] (0,0) -- (0,3.2);
    \draw[blue!70,thick,domain=0:1,samples=60,variable=\x]
      plot ({\x*4.2},{2.7*\x*\x});
    \draw[orange!80!black,thick,dashed,domain=0:1,samples=60,variable=\x]
      plot ({\x*4.2},{2.7*(1-\x)*(1-\x)});
    \node[font=\footnotesize,black!75,anchor=west] at (2.5,2.55) {校准偏差};
    \node[font=\footnotesize,black!75,anchor=west] at (0.15,2.55) {假阳率差};
    \node[font=\footnotesize,anchor=north] at (2.1,-0.25) {权衡参数};
    \node[font=\footnotesize,anchor=south] at (2.1,3.25) {基率不同};
    \draw[black!40,dotted] (2.1,0) -- (2.1,0.675);
    \node[font=\footnotesize,black!65,anchor=north] at (2.1,-0.85) {没有公共零点};
  \end{scope}
  % ---- 右图：基率相同 ----
  \begin{scope}[xshift=6.9cm]
    \draw[->,black!60] (0,0) -- (4.7,0);
    \draw[->,black!60] (0,0) -- (0,3.2);
    \draw[blue!70,thick,domain=0:1,samples=60,variable=\x]
      plot ({\x*4.2},{2.7*(2*\x-1)*(2*\x-1)});
    \draw[orange!80!black,thick,dashed,domain=0:1,samples=60,variable=\x]
      plot ({\x*4.2},{1.6*(2*\x-1)*(2*\x-1)});
    \node[font=\footnotesize,anchor=north] at (2.1,-0.25) {权衡参数};
    \node[font=\footnotesize,anchor=south] at (2.1,3.25) {基率相同（对照）};
    \fill[black!70] (2.1,0) circle (0.075);
    \node[font=\footnotesize,black!65,anchor=north] at (2.1,-0.85) {两者同时触底};
  \end{scope}
\end{tikzpicture}

\emph{左图两条曲线没有公共零点，这是命题的图形形态；右图把基率差撤掉，冲突随之消失。决定有没有取舍的是基率差，不是模型有多好。}
\end{center}

左图那条虚线在横轴左端触底，实线在右端触底，而它们永远不同时落地。上一节那两份报告就分别站在这条横轴的两端：开发方站在左端，批评方站在右端，各自看到自己那条曲线为零，也各自看不到对方那条正在升高。右图是唯一能让争论自然消解的情形，而它要求的前提——两个群体的 \(Y\) 分布相同——不是模型能提供的。

\subsection{这条命题没有说的事}

它最常被读成``公平不可能''。这个读法是错的，而且错得有后果。

命题的实际内容是：\textbf{``公平''这个词必须先被指定成一个具体判据，指定之后才谈得上做到与否}。在指定之前争论，双方各握一个真数字，会永远各说各话——上一节那两份报告就是这个状态的标准形态。命题做的事是把``要哪一个''从一个可以回避的问题变成一个必须当场回答的问题。

指定判据这件事本身不是数学问题。哪一类错误在这个场景里更该被压低，取决于误判的后果落在谁身上、以及这个系统的判定进入了怎样的下游流程——这是\hyperref[sec:11-Mathematical-Statistics/07-Statistical-Decision-and-Reliable-Prediction/01]{``状态、行动、损失与风险构成一个决策问题''}里那套语言要处理的：先写下损失矩阵，判据的选择才有依据。同样地，\hyperref[sec:18-Synthesis/01-System-Lifecycle/01]{``把诉求写成目标函数与约束''}讲的``把诉求翻译成目标''，在这里就是把某个判据写成硬约束、把另一个写进目标——两者不能都是硬约束，可行域会空。

还有一件事命题没有说：它没有断言在给定基率差下，各判据之间只能沿一条线取舍。可行的 \((\mathrm{TPR},\mathrm{FPR})\) 组合构成一个二维区域，落在边界上的方案互不支配，选哪一个不由数学决定。真正被排除掉的只有一个点——那个让所有判据同时归零的点。

上面所有判据有一个共同点还没被检查：它们全都只用到联合分布 \(P(A,S,Y)\)。这意味着两个内部机制完全不同的世界，只要产生同样的数据，在这些判据下就一模一样。下一节把这一点推到底，看看当人们真正想问的是``换个群体这个决定会不会变''时，观察层面的判据够不够用。
